\documentclass[10pt,a4paper]{amsart}

\pdfinfoomitdate=1
\pdftrailerid{}
\pdfsuppressptexinfo=15

\usepackage[T1]{fontenc}
\usepackage{lmodern}
\usepackage[margin=1in]{geometry}
\usepackage{microtype}
\usepackage{amsmath,amssymb,mathtools}
\usepackage{booktabs}
\usepackage{natbib}
\usepackage[hidelinks]{hyperref}

\hypersetup{
  pdftitle={},
  pdfauthor={},
  pdfsubject={},
  pdfkeywords={}
}

\newcommand{\cC}{\mathcal C}
\newcommand{\cP}{\mathcal P}
\newcommand{\EE}{\mathbb E}
\newcommand{\PP}{\mathbb P}
\newcommand{\one}{\mathbf 1}
\newcommand{\Lam}{\Lambda}
\newcommand{\set}[1]{\left\{#1\right\}}
\DeclareMathOperator{\Var}{Var}
\newtheorem{theorem}{Theorem}[section]
\newtheorem{corollary}[theorem]{Corollary}

\title{Recorded-Transversal Laws on Rate-Weighted Sunflower Forests}
\author{Anonymous}
\date{}

\begin{document}

\begin{abstract}
Consider the owned random-edge/random-vertex covering process, restricted to
vertex-disjoint heterogeneous sunflowers with fixed positive edge rates.  We
derive the complete law of the full recorded transversal, without reducing it
to a minimal transversal.  For one sunflower, an inclusion--exclusion integral
gives every selected-petal mask at arbitrary rates, and independent uniform
marks refine that law to every actual vertex set.  At unit rates, elementary
symmetric polynomials give the complete selection-count probability generating
function, mean, and variance; the maximal-count atom is resolved into its
all-petal and final-core events.  For a disjoint sunflower forest, the marked
stopped laws factor; component choice counts add, hence their probability
generating functions multiply.  The covering algorithm,
sunflower carrier, exponential/size-biased ordering, and independence of
disjoint restrictions are prior or generic machinery and are not claimed as
contributions.  A finite exact-arithmetic program checks the formulas on a
frozen bounded family.
\end{abstract}

\maketitle

\section{Positioning and contribution}

Sunflowers, also called Delta-systems, are classical
\citep{erdosrado1960}.  Their transversals belong to a broad literature that
includes extremal and approximation questions \citep{chvatalmcdiarmid1992}.
More importantly for the present dynamics, the random covering algorithm is
not new.  Pitt gave its graph ancestor \citep{pitt1985}, and
\citet[Section~5.1 and Theorem~6]{baryehuda2000} directly treats the Pitt rule
for hypergraph vertex cover.  In the unweighted setting, its vertex choice in
the selected hyperedge is uniform.  We therefore assign zero credit to the
covering algorithm, its validity, and every approximation guarantee.

The ordering device is also owned background.  Rate-proportional sampling
without replacement is a Plackett--Luce/size-biased order
\citep{plackett1975}; positive-rate independent exponentials and dissociation
of restrictions to disjoint item sets are treated explicitly by
\citet{gnedin2026}.  We assign zero credit to the exponential representation,
memorylessness, finite ranking probabilities, and generic forest
independence.  The same applies to inclusion--exclusion, beta integrals,
elementary symmetric polynomials, tail-sum identities, and products of PGFs.

After this subtraction, the contribution is deliberately narrow: for the
owned covering process restricted to vertex-disjoint heterogeneous sunflowers
with fixed edge rates, we derive the complete recorded-transversal law and, at
unit rates, the complete choice-count PGF and first two moments.  The point is the
closed conjunction of aggregate endpoint, actual-vertex refinement,
choice-count law, and marked forest factorization for this carrier.  We claim neither a new
algorithm nor a general theory of random greedy hitting sets.

\section{Model and marked order}

Fix positive integers
\[
  m,c,p_1,\ldots,p_m
\]
and positive real numbers $\lambda_1,\ldots,\lambda_m$.  A sunflower has a
core $C$ of size $c$, pairwise disjoint petals $P_i$ of size $p_i$, and edges
\[
  E_i=C\mathbin{\dot\cup}P_i,\qquad i\in[m].
\]
At each step, an active edge $E_i$ is chosen with probability proportional to
its fixed rate $\lambda_i$.  A vertex is then chosen uniformly from $E_i$ and
added to the recorded set.  A petal choice deletes only $E_i$; a core choice
hits all remaining edges and stops the component.  If every edge receives a
petal choice, the process also stops.  The endpoint is the full recorded set:
previously selected petal vertices are retained, so the endpoint need not be a
minimal transversal.

Write
\[
 r_i=\frac{p_i}{c+p_i},\qquad
 q_i=\frac{c}{c+p_i}=1-r_i,
 \qquad
 \Lam(A)=\sum_{i\in A}\lambda_i.
\]
Let $X_i\sim\mathrm{Exp}(\lambda_i)$ independently.  Independently of all
clocks and one another, give edge $i$ a uniform mark $U_i\in E_i$.  Sorting
the $X_i$ and reading marks until the first core mark gives the original
process.  Indeed, the minimum clock has the required rate-proportional law,
and after a petal deletion the residual clocks have the same law by
memorylessness.  The fixed-rate and independent-mark hypotheses are essential
for this coupling.  This representation is used only as owned machinery.

\section{The complete weighted endpoint law}

For a proper subset $A\subsetneq[m]$, let $\pi(A)$ be the probability that
exactly the edges in $A$ are recorded through petal vertices before a core
vertex stops the component.  Define
\begin{equation}\label{eq:integral-sum}
 I(A)=\sum_{B\subseteq A}
       \frac{(-1)^{|B|}}{\Lam([m]\setminus A)+\Lam(B)}.
\end{equation}
Every denominator is positive because $A$ is proper.

\begin{theorem}[weighted aggregate endpoint]\label{thm:weighted}
For every proper $A\subsetneq[m]$,
\begin{equation}\label{eq:weighted-endpoint}
 \pi(A)=
 \left(\prod_{i\in A}r_i\right)
 \left(\sum_{j\notin A}q_j\lambda_j\right)I(A).
\end{equation}
The remaining, all-petal endpoint has probability
\begin{equation}\label{eq:all-petal}
 \pi([m])=\prod_{i=1}^{m}r_i.
\end{equation}
These $2^m$ masses are positive and sum to one.
\end{theorem}

\begin{proof}
Fix $A\subsetneq[m]$ and expose the time $t$ of the stopping core mark.  For
each $i\in A$, edge $i$ must carry a petal mark and satisfy $X_i<t$.  For
each edge outside $A$, no clock may ring before $t$, except that one such edge
$j$ rings at $t$ and carries a core mark.  Summing the latter density over
$j\notin A$ gives
\begin{align}
 \pi(A)
 &=\left(\prod_{i\in A}r_i\right)
   \left(\sum_{j\notin A}q_j\lambda_j\right)
   \int_0^\infty
   e^{-\Lam([m]\setminus A)t}
   \prod_{i\in A}(1-e^{-\lambda_i t})\,dt.\label{eq:weighted-integral}
\end{align}
Expanding the product in \eqref{eq:weighted-integral} and integrating each
exponential gives \eqref{eq:integral-sum} and
\eqref{eq:weighted-endpoint}.  The only way to terminate without a core mark
is for every mark to lie in its petal, proving \eqref{eq:all-petal}.
The marked-clock construction partitions its probability space into these
events, which proves normalization; the integral form proves positivity.
\end{proof}

The aggregate law also resolves every actual vertex set, including at
unequal rates.

\begin{corollary}[actual-vertex law]\label{cor:vertices}
Fix $A\subsetneq[m]$, vertices $x_i\in P_i$ for $i\in A$, and $y\in C$.  Then
\begin{equation}\label{eq:actual-core}
 \PP\!\left(
  R=\set{y}\cup\set{x_i:i\in A}
 \right)
 =\frac{\pi(A)}{c\prod_{i\in A}p_i}.
\end{equation}
For $x_i\in P_i$ for every $i$,
\begin{equation}\label{eq:actual-all-petal}
 \PP\!\left(R=\set{x_1,\ldots,x_m}\right)
 =\prod_{i=1}^{m}\frac{1}{c+p_i}.
\end{equation}
\end{corollary}

\begin{proof}
Conditional on its category, every petal mark is uniform on its own $P_i$
and the stopping core mark is uniform on $C$.  More formally, let
$\mathcal F$ be the sigma-field generated by all clocks and all core-versus-
petal category indicators.  The event fixing $A$ and the terminal edge is
$\mathcal F$-measurable; conditional on $\mathcal F$, the relevant
within-category vertex identities remain mutually independent and uniform on
$C$ or their respective petals.  Dividing the aggregate mass by
$c\prod_{i\in A}p_i$ proves \eqref{eq:actual-core}.  In the all-petal event,
edge $i$ selects a specified petal vertex with probability $1/(c+p_i)$;
independence gives \eqref{eq:actual-all-petal}.
\end{proof}

\section{Unit rates: choice-count PGF and moments}

Assume throughout this section that $\lambda_i=1$.  Let $T\in\set{1,\ldots,m}$
be the discrete number of selections, equivalently recorded vertices, at
absorption.  Thus $T$ is a choice count, not elapsed time in the exponential
embedding.  Let $e_t(r_1,\ldots,r_m)$ be the elementary symmetric polynomial,
and put
\begin{equation}\label{eq:st}
 s_t=\frac{e_t(r_1,\ldots,r_m)}{\binom mt},\qquad 0\leq t\leq m,
\end{equation}
with $s_0=1$.

\begin{theorem}[complete choice-count law]\label{thm:stopping}
For $0\leq t<m$,
\begin{equation}\label{eq:tail}
 \PP(T>t)=s_t.
\end{equation}
Consequently,
\begin{equation}\label{eq:mass}
 \PP(T=t)=s_{t-1}-s_t\quad(1\leq t<m),
 \qquad
 \PP(T=m)=s_{m-1}.
\end{equation}
The maximal-count atom is the disjoint union of two mechanisms:
\begin{equation}\label{eq:top-atom}
 \PP(T=m)
 =\prod_{i=1}^m r_i
  +\frac1m\sum_{j=1}^m q_j\prod_{i\ne j}r_i
 =\frac{e_{m-1}(r_1,\ldots,r_m)}{m}.
\end{equation}
In particular, its first term is the all-petal endpoint, while its second term
is a final core choice after $m-1$ petal choices.
\end{theorem}

\begin{proof}
At unit rates the clock order is a uniform random permutation of $[m]$,
independent of the marks.  The event $T>t$ says that the first $t$ distinct
edges in this order all have petal marks.  Averaging
$\prod_{i\in A}r_i$ over the $\binom mt$ possible first-edge sets gives
\eqref{eq:tail}, and differencing tails gives \eqref{eq:mass} below the top
atom.  For $T=m$, either all $m$ marks are petals, or a unique edge $j$ is
last in the uniform order, its mark is core, and all other marks are petals.
This gives the first equality in \eqref{eq:top-atom}.  Substituting
$q_j=1-r_j$ reduces it to
$m^{-1}\sum_j\prod_{i\ne j}r_i=e_{m-1}/m=s_{m-1}$.
\end{proof}

\begin{corollary}[PGF, mean, and variance]\label{cor:moments}
The probability generating function and first two moments are
\begin{align}
 G(z)&=\sum_{t=1}^{m-1}(s_{t-1}-s_t)z^t+s_{m-1}z^m,
 \label{eq:pgf}\\
 \EE T&=\sum_{t=0}^{m-1}s_t,\label{eq:mean}\\
 \EE T^2&=\sum_{t=0}^{m-1}(2t+1)s_t,\label{eq:second}\\
 \Var(T)&=\sum_{t=0}^{m-1}(2t+1)s_t-
 \left(\sum_{t=0}^{m-1}s_t\right)^2.\label{eq:variance}
\end{align}
\end{corollary}

\begin{proof}
Equation \eqref{eq:pgf} is the mass formula
\eqref{eq:mass}.  For every positive integer-valued $T\leq m$,
$T=\sum_{t=0}^{m-1}\one_{\{T>t\}}$ and
$T^2=\sum_{t=0}^{m-1}(2t+1)\one_{\{T>t\}}$.  Taking expectations and using
\eqref{eq:tail} proves \eqref{eq:mean}--\eqref{eq:variance}.
\end{proof}

The heterogeneity is visible in every $e_t$: replacing the $r_i$ by their
mean generally changes all terms of degree at least two and does not preserve
the endpoint law.

\section{Disjoint sunflower forests}

Let $H_1,\ldots,H_d$ be vertex-disjoint sunflower components, each with its
own positive parameters and fixed edge rates.  The global scheduler chooses
among all active edges with probability proportional to their rates.  Write
$R_a,T_a$, and $G_a$ for the recorded endpoint, discrete choice count, and
choice-count PGF in component $a$; write $R$ and $T$ for the forest
quantities.

\begin{theorem}[marked stopped factorization]\label{thm:forest}
For arbitrary positive fixed rates,
\begin{equation}\label{eq:forest}
 \mathcal L(R)=\bigotimes_{a=1}^d\mathcal L(R_a),
 \qquad
 T=\sum_{a=1}^dT_a,
 \qquad
 G_T(z)=\prod_{a=1}^dG_a(z).
\end{equation}
Thus the endpoint mass of a forest endpoint is the product of the component
masses from Theorem~\ref{thm:weighted} and
Corollary~\ref{cor:vertices}; at unit rates each $G_a$ is given by
\eqref{eq:pgf}.
\end{theorem}

\begin{proof}
Assign independent clocks and marks to every edge.  Each local stopped marked
history $(R_a,T_a)$ is a function only of the clocks and marks in $H_a$;
hence these pairs are independent.  Sorting all clock rings across components
gives the global rate-proportional scheduler.  Once a core mark stops one
component, suppressing its later clocks does not change any local history.
The forest records exactly the union of the local records and makes exactly
the sum of their choices.  Tensorization and multiplication of PGFs follow.
\end{proof}

The theorem's content is only the fully marked stopped consequence for this
process.  Independence/dissociation of disjoint exponential-order
restrictions is generic and receives no novelty credit.

Continuous elapsed absorption time is not analyzed here.  If one equips the
forest with the exponential embedding above and writes $S_a$ for component
$a$'s continuous stopping time, then the whole forest completes at
$S=\max_a S_a$, not at $\sum_a S_a$; no wall-clock convolution is claimed.

\section{Exact controls and limitations}

The accompanying self-contained Python program uses only integers and
\texttt{fractions.Fraction}.  It performs no sampling or floating-point
calculation.  It recursively enumerates the process and compares every mass
with the displayed formulas.  The frozen scope is:
\begin{center}
\begin{tabular}{@{}p{0.24\textwidth}p{0.62\textwidth}r@{}}
\toprule
lane & parameter range & inputs\\
\midrule
unit-rate aggregate & $c\in\{1,2,3\}$, $m\in\{1,\ldots,5\}$,
  $p_i\in\{1,\ldots,4\}$ & 4092\\
weighted aggregate & $c\in\{1,2\}$, $m\in\{1,2,3\}$,
  $p_i,\lambda_i\in\{1,2,3\}$ & 1638\\
unit-rate actual vertices & $c\in\{1,2\}$, $m\in\{1,2,3\}$,
  $p_i\in\{1,2,3\}$ & 78\\
two-component forests & three unit-rate and one weighted control & 4\\
\bottomrule
\end{tabular}
\end{center}
These are falsification controls, not proofs.  In particular, the program does
not exhaust weighted vertex-resolved endpoints or arbitrary forests; those
claims follow from the all-parameter arguments above.  The bounded owner
search located the direct process and ordering owners already subtracted, but
did not locate this complete conjunction.  That non-hit is not a novelty
certificate, and a direct package owner would supersede the present
positioning.

\section{Conclusion}

On a heterogeneous sunflower, the marked size-biased order separates the
recorded-transversal problem into a weighted endpoint integral, independent
actual marks, and---at unit rates---an elementary-symmetric choice-count law.  On
vertex-disjoint components, these marked stopped laws tensorize.  The result
is a complete exact atlas for a sharply restricted owned process, not a new
covering algorithm, ordering method, or independence principle.

\bibliographystyle{plainnat}
\bibliography{references}

\end{document}
